% Do not edit by hand; generated from paper/datasets.Rmd.
This appendix describes the 16 corpora used in the paper. Note that the
full BEAR database contains additional datasets. More complete
documentation is available online.

\newpage

\subsection{psymetadata and Nuijten et al (psychology
datasets)}\label{psymetadata-and-nuijten-et-al-psychology-datasets}

\textbf{Reference:} A set of psychology datasets compiled by
\textcite{rodriguez2022psymetadata}. We also disaggregate a large
meta-meta-analysis of intelligence research
\autocite{nuijten2020intelligence}.

\textbf{Snapshot date:} January 2026

\textbf{Research question:} \emph{psymetadata} is an R package that
curates multiple open datasets from published meta-analyses in
psychology for teaching and demonstrations.
\textcite{nuijten2020intelligence} is a large meta-meta-analysis of
intelligence research, which focused on power and small study effects.

\textbf{Data availability:} 22 datasets are distributed with the
\texttt{psymetadata} R package (CRAN). We do not use its Many Labs 2
dataset, which we obtain from the original source and include
separately. After excluding Many Labs 2 and separately representing the
intelligence meta-meta-analysis, the remaining psychology datasets range
from 65 rows to over 1,000. The package is GPL-3; individual datasets
originate from the cited source papers.

\textbf{Data processing:} Processing was minimal. We used regular
expressions to extract year labels from titles. We dropped some invalid
rows (missing, non-finite, or non-positive standard errors). Signed
z-values were computed as effect sizes divided by their standard errors.

\textbf{Study characteristics:}

\begin{itemize}
\item
  \textbf{psymetadata:} the designs of the primary studies are not
  harmonised across the imported datasets.
\item
  \textbf{psymetadata:} the effect-size measure is not harmonised across
  the imported datasets.
\item
  \textbf{psymetadata:} records are grouped by original psymetadata
  dataset.
\item
  \textbf{Nuijten:} the designs of the primary studies are not
  classified.
\item
  \textbf{Nuijten:} BEAR retains the reported effect estimates and
  standard errors without a harmonised effect-size measure.
\end{itemize}

\newpage

\subsection{Askarov et al 2023
(economics)}\label{askarov-et-al-2023-economics}

\textbf{Reference:} \textcite{askarov2023significance}.

\textbf{Research question:} Impact of mandatory data sharing on
(excessive) statistical significance in economics papers

\textbf{Data availability:} The file
\texttt{Mandatory\ data-sharing\ 30\ Aug\ 2022.dta} may be downloaded
from
\url{https://github.com/anthonydouc/Datasharing/blob/master/Stata/}; the
reproducibility package is available at
\url{https://dataverse.harvard.edu/file.xhtml?fileId=7884702&version=1.0}.

\textbf{Data description and source:} The paper includes 64,000 economic
estimates surveyed by Ioannidis et al.~(2017) plus results of
``searching through numerous databases, journal websites, and also an
email survey of 109 authors known to have produced meta-analyses of
economic topics. We cast the widest net possible to capture empirical
economic estimates. Thus, some of the papers included in these
meta-analyses were originally published in political science and
psychology journals. Nevertheless, these meta-studies deal with
economics topics.'' Later they note: ``we collect 166,924 parameter
estimates from 345 distinct research areas, of which 20,121 were
published in 24 of the leading general interest and field journals.''
The paper mentions 12,521 observations from before and 2,537 from after
the introduction of data-sharing policies in 24 journals; the whole
reproducibility dataset has 32 journals (the authors removed some
smaller ones) and 22,172 rows.

\textbf{Notes:} The article by Askarov et al.~starts from the
Ioannidis-Stanley-Doucouliagos dataset, which is also used by
Arel-Bundock et al.~(another dataset in BEAR). Both datasets incorporate
additional data and target different fields of research, but there is
likely substantial overlap between the two sets.

\textbf{Data processing:} We use the same source file
\texttt{Mandatory\ data-sharing\ 30\ Aug\ 2022.dta} (11.6 MB) as in the
reproducibility package. Signed z-values were computed as reported
effect sizes divided by their standard errors.

\textbf{Study characteristics:}

\begin{itemize}
\item
  \textbf{Study ID:} we use the source study identifier.
\item
  \textbf{Meta-analysis ID:} we use \texttt{filename}.
\item
  \textbf{Topic:} we use the source's flag to distinguish macroeconomics
  or microeconomics/other.
\item
  \textbf{Method:} The source flags experimental and mixed
  observational/experimental research. We label estimates with neither
  flag set as ``not RCT, not mixed''. This category may include
  quasi-experimental research, which the source does not separately
  identify.
\item
  We do not have access to the measure of effect size (e.g.~ratio vs
  mean difference).
\end{itemize}

\newpage

\subsection{Metapsy (psychotherapy)}\label{metapsy-psychotherapy}

Studies included in \emph{Metapsy} (\url{https://www.metapsy.org/}), a
set of living meta-analytic databases of randomised trials of
psychotherapeutic interventions.

\textbf{Snapshot date:} January 2026

\textbf{Data availability:} Metapsy is maintained by an international
collaboration led by Vrije Universiteit Amsterdam. Databases are
available at \url{https://www.metapsy.org/} and through the R package
\texttt{metapsyData} \autocite{harrer2022metapsydata}. The Metapsy FAQ
states that all data provided as part of the project are open access and
that attribution is required. The source contains RCTs of
psychotherapeutic interventions.

\textbf{Data processing:} We merged 20 individual datasets in Metapsy
and performed some cleaning (e.g.~to extract study dates from titles).
Almost all estimates are Hedges' g, but for one of the databases
(\texttt{total-response}) we converted log-odds ratios to SMD (dividing
by \(\Pi/\sqrt{3}\) to scale SD of logistic regression) to keep a
comparable scale. Signed z-values were computed as effect sizes divided
by their standard errors.

\textbf{Study characteristics:}

\begin{itemize}
\item
  \textbf{Sample size and year:} \texttt{ss} is the sum of
  \texttt{n\_arm1} and \texttt{n\_arm2} when both are supplied;
  \texttt{year} is the source year or a year parsed from the study
  label.
\item
  \textbf{Study ID:} we use the source \texttt{study} label within each
  database.
\item
  \textbf{Meta-analysis ID:} we use the Metapsy database identifier.
\item
  \textbf{Topic:} we use an explicit clinical-subject lookup; datasets
  \texttt{peer-support} and \texttt{total-response} are labelled as
  ``mental health''.
\item
  Effects are standardised mean differences, mainly Hedges' g; a
  minority are converted from log odds ratios.
\end{itemize}

\newpage

\subsection{Arel-Bundock et al 2022 (political
science)}\label{arel-bundock-et-al-2022-political-science}

\textbf{Reference:} \textcite{arel2022quantitative}.

\textbf{Research question:} assess statistical power in political
science research

\textbf{Data availability:} The replication package (a zip file dated
\texttt{20241010}) is available at \url{https://osf.io/fgdet}. All data
are in the public domain.

\textbf{Data description and source:} The authors searched for
meta-analytic articles across 141 journals in political science and
contacted their authors. The resulting source dataset contains 16,649
hypothesis tests, grouped in 351 meta-analyses and reported in 46
peer-reviewed meta-analytic articles.

\textbf{Notes:} This dataset collects meta-analysis inputs, rather than
repeating data extraction from individual studies. It combines an
earlier dataset from Doucouliagos et al with another search. Therefore,
some of this dataset will overlap with Askarov et al. However, the two
articles aim to cover different fields and BEAR keeps the two datasets
separate.

\textbf{Data processing:} We used the same code as the original authors
(a replication package that uses makefiles), but skipped some additional
processing they performed for their paper. BEAR uses
\texttt{question\_id} as \texttt{metaid} (351 groups defined by
meta-analysis or research question). Signed z-values were computed as
reported estimates divided by their standard errors. We verified that
the distributions of z-values matched those in the article.

\textbf{Study characteristics:}

\begin{itemize}
\item
  \textbf{Study ID:} we use author's study ID; for \textbf{meta-analysis
  ID} we use \texttt{question\_id}.
\item
  \texttt{subset} column is the source \texttt{meta\_id}
\item
  We do not have information on the study design or measure of effect
  size.
\end{itemize}

\newpage

\subsection{Costello and Fox 2022; Yang et al 2023, 2024 (ecology and
evolution)}\label{costello-and-fox-2022-yang-et-al-2023-2024-ecology-and-evolution}

\textbf{Reference:} These two datasets come from the same replication
package and were processed using the shared pipeline in Yang et al.~2024
\autocite{yang2024large}. The original references for the datasets are
Yang et al.~2023 and Costello and Fox 2022
\autocite{yang2023publication,costello2022decline}.

\textbf{Research question:} decline effects in ecology (Costello and
Fox) and publication bias and performance of empirical research (Yang et
al)

\textbf{Data collection:} Costello and Fox examined 466 meta-analyses
identified by searching Web of Science. The reproducibility dataset
includes 88,000 rows across 466 analyses in 232 papers. Subsequently,
Yang et al.~removed duplicated studies and zeroes, discarding about 20\%
of rows. Yang et al.~collected 102 meta-analyses published in relevant
journals between 2010 and 2019, after removing records duplicated in the
Costello and Fox dataset.

\textbf{Data availability:} The two datasets may be downloaded from
\url{https://github.com/Yefeng0920/replication_EcoEvo_git/}. They are
available under a CC BY 4.0 licence.

\textbf{Data processing:} The only additional processing was to clean up
study years. Signed z-values were computed as effect sizes divided by
their standard errors. For Yang et al, BEAR uses the source file names
as meta-analysis identifiers.

\textbf{Costello and Fox:}

In \texttt{BEAR.rds}, we re-label the effect measures as SMD, response
ratio and correlation. The individual dataset keeps the original Yang et
al.~coding (Zr, lnRR, SMD).

\textbf{Yang:}

In \texttt{BEAR.rds}, we code Fisher's z as ``correlation'' and lnRR as
``response\_ratio'' for consistency with other datasets.

\textbf{Study characteristics:}

\begin{itemize}
\item
  \textbf{Study ID:} for Costello and Fox, we use \texttt{study2}; for
  Yang, we use \texttt{study\_ID}.
\item
  \textbf{Meta-analysis ID:} for Costello and Fox, we use
  \texttt{meta.analysis.id}; for Yang, we use \texttt{meta\_id}.
\item
  \textbf{Topic:} for Costello and Fox, we use ecology, evolution or
  other ecology/evolution. Yang has no topic classification.
\item
  \textbf{Costello and Fox:} the primary-study design is not classified.
\item
  \textbf{Costello and Fox:} effects are standardised mean differences,
  log response ratios, Fisher's z correlations, or less common measures.
  Less common measures include regression coefficients, log odds ratios
  and mean differences. Some standardised mean differences are recorded
  as absolute values.
\item
  \textbf{Yang:} the primary-study design is not classified.
\item
  \textbf{Yang:} effects are standardised mean differences, log response
  ratios, or Fisher's z correlations.
\end{itemize}

\newpage

\subsection{What Works Clearinghouse
(education)}\label{what-works-clearinghouse-education}

Large database of education effect sizes created by the Institute of
Education Sciences \autocite{wwc2025studyfindings}.

\textbf{Snapshot date:} 17 May 2025

\textbf{Data availability:} Flat files are available at
\url{https://ies.ed.gov/ncee/wwc/studyfindings}. The Permissions and
Disclaimers section of that website states: ``Unless stated otherwise,
all information on the U.S. Department of Education's IES website
\ldots{} is in the public domain and may be reproduced, published,
linked to, or otherwise used without permission from IES.''

\textbf{Data collection:} The WWC source data have 13,054 findings (data
rows) across 1,908 reviews. Each row is one finding, so a single source
paper can contribute multiple rows. The dataset includes effect sizes,
but not standard errors, as well as p-values reported by the studies and
p-values re-calculated by the WWC reviewers. After BEAR processing and
filtering, the main BEAR dataset includes 7,080 findings.

\textbf{Data processing:} We include only RCTs and quasi-experimental
studies (98\% of data altogether). Majority of these (about 80\%) are
RCTs.

We removed a few rows where the results were flagged as not meeting WWC
standards or as ineligible for review. We also removed rows marked by
WWC as subgroup measures to avoid treating subgroup analyses as
independent focal findings. Where available, we preferred effect sizes
and p-values calculated by WWC to the ones reported by the studies. We
replaced p-values that were equal to zero (this occurred in 3\% of the
dataset) with a truncated p-value (i.e.~assumed \(p < 10^{-16}\)). We
calculated z-values from these p-values under the assumption of
two-sided tests. We then applied the sign of the reported effect size,
so z-values are signed where the effect size is non-zero.

\textbf{Study characteristics:}

\begin{itemize}
\item
  \textbf{Study ID:} we use citation-derived author and year; we use the
  original \texttt{WWCStudy} as fall-back in some cases.
\item
  \textbf{Topic:} WWC dataset includes some useful topic
  classifications, but these are not included in \texttt{BEAR.rds} as
  they are either too granular or are overlapping flags for study
  categories
\end{itemize}

\newpage

\subsection{Cochrane Database of Systematic Reviews
(CDSR)}\label{cochrane-database-of-systematic-reviews-cdsr}

Trials with results available in CDSR,
\url{https://www.cochranelibrary.com/} \autocite{cochrane2025cdsr}.

\textbf{Snapshot date:} 20 November 2025

\textbf{Data collection:} Using the R package \texttt{cochrane}
\autocite{schwab2024cochrane}, we combined data from 8,726 CDSR review
records into a single file, \texttt{cdsr\_interventions\_19nov2025.csv}
(63 MB).

\textbf{Data processing:} We applied several processing steps and
extensive filtering.

During processing, we cleaned up study years and categorised measures
(risk ratio, odds ratio, mean difference, etc.) and experimental
designs. In particular, we created an ``RCT'' flag by scanning each
review abstract for its inclusion criteria. Most studies in CDSR are
RCTs, but some reviews also allowed quasi-experimental studies, so we
could not classify many studies. We also classify outcome rows into
\texttt{efficacy}, \texttt{safety}, \texttt{dropouts}, and \texttt{bias}
based on the comparison, outcome, and subgroup labels.

To calculate effect sizes and z-values for BEAR, we used \emph{metafor}
to calculate standardised mean differences for continuous outcomes and
probit-transformed differences for dichotomous outcomes. This puts
continuous and binary outcomes on a comparable scale.

Unlike in most of the other datasets in BEAR, we filter CDSR data
heavily. There are five important decisions:

\begin{enumerate}
\def\labelenumi{\arabic{enumi}.}
\item
  Use only the outcome and comparison from each review coded as ``1'',
  as these are most likely to be the primary outcome and most relevant
  comparison.
\item
  Keep only rows classified as \texttt{efficacy}, excluding likely
  safety, dropout, and bias-related outcomes based on comparison,
  outcome, and subgroup labels.
\item
  Use only studies with continuous and dichotomous outcomes, removing
  rows where effect estimates are based on instrumental variables or
  individual patient data.
\item
  Remove rows where the effect measure is unknown (keeping OR, RR, Peto
  OR, mean difference, standardised mean difference, and risk
  difference) and rows with no participants.
\item
  Remove count data with no events or a 100\% event rate in both arms
  (see the discussion of zero-event studies).
\end{enumerate}

These processing steps are applied in the main dataset,
\texttt{BEAR.rds}. Additional rows are still available to researchers in
the individual dataset, \texttt{Cochrane.rds}.

As a result, the distributed \texttt{BEAR.rds} file includes a smaller
analysis subset with about 36,000 rows: the first comparison and first
outcome from each review, restricted to efficacy outcomes. The processed
\texttt{data/Cochrane.rds} file is larger and is available for users who
want to make different selection decisions. It contains 760,486 result
rows, representing 90,042 review-study pairs across 6,619 Cochrane
reviews.

A large proportion of CDSR rows have binary outcomes with no events in
either arm or events in all participants in both arms. In the BEAR
analysis subset before this exclusion, 2,371 of 26,619 binary rows
(8.9\%) met one of these conditions: 2,247 rows (8.4\%) had no events in
either arm and 124 rows (0.5\%) had all participants experiencing the
event in both arms.

Zero-event studies may be important for some metascientific analyses,
but for modelling the distribution of z-values they are unhelpful. With
these rows included, the empirical distribution has extra observations
at or near zero, and the normal approximation is poor. Because of this
peak at zero, including these non-significant studies has the
counter-intuitive effect of increasing the estimated degree of
selection: the fitted distribution has a steeper drop-off between zero
and the conventional significance threshold at 1.96.

\textbf{Study characteristics:}

\begin{itemize}
\item
  \textbf{Sample size and year:} \texttt{ss} is the sum of the two arm
  totals; \texttt{year} is \texttt{study.year}.
\item
  \textbf{Study ID:} we use \texttt{study.name}.
\item
  \textbf{Meta-analysis ID:} we use the Cochrane review identifier
  (\texttt{id}).
\item
  \textbf{Topic:} we use the Cochrane specialty.
\item
  \texttt{subset} records study-data provenance: published, unpublished,
  sought or mixed.
\item
  \texttt{method}: We review meta-analysis eligibility criteria to
  determine which meta-analyses included only randomised trials. This
  classifies 69\% of estimates as RCTs; the remaining 31\% could not be
  classified reliably, although many are likely to be RCTs because
  Cochrane meta-analyses typically include RCT or quasi-randomised
  evidence.
\item
  \texttt{measure}: About 32\% of estimates are standardised mean
  differences and 68\% are differences between probit-transformed event
  proportions.
\item
  \texttt{topic} is Cochrane specialty and \texttt{subset} is source
  data type (published, unpublished, sought, or mixed);
  \texttt{outcome\_group} is also available in the processed data.
\end{itemize}

\newpage

\subsection{Chavalarias et al.~2016
(PubMed/Medline)}\label{chavalarias-et-al.-2016-pubmedmedline}

\textbf{Reference:} \textcite{chavalarias2016evolution}.

\textbf{Research question:} reporting of p-values in the biomedical
literature

\textbf{Data collection:} The authors used large-scale text mining of
MEDLINE abstracts and PubMed full texts: 4.5 million p-values in 1.6
million MEDLINE abstracts, and 3.4 million p-values in 385,000 PubMed
full-text articles.

\textbf{Data availability:} The extracted p-value dataset is publicly
available at
\url{https://dataverse.harvard.edu/dataset.xhtml?persistentId=doi:10.7910/DVN/6FMTT3}
(about 6 GB of data). According to the hosting platform, this dataset is
released under the Creative Commons Attribution-NonCommercial 4.0
International Licence.

\textbf{Data processing:} Unsigned z-values were derived from p-values
assuming two-sided tests. P-value truncation operators were processed by
collapsing variants such as \texttt{\textless{}\textless{}},
\texttt{\textless{}\textless{}\textless{}}, \texttt{\textless{}=},
\texttt{less\ than}, and \texttt{=\textless{}} into
\texttt{\textless{}}. We dropped 0.7\% of rows where p-values did not
have a ``plain'' format and 0.08\% of rows where truncation could not be
unambiguously classified. For the distributed version of BEAR 50,000
studies are selected at random, with one row per study, in order to keep
the file size manageable.

\textbf{Study characteristics:}

\begin{itemize}
\item
  \textbf{Study ID:} we use the PubMed identifier (\texttt{PMID}).
\item
  \texttt{subset} distinguishes estimates obtained by authors from
  MEDLINE abstracts vs from PMC full text.
\item
  There is no information on study design (i.e.~we cannot distinguish,
  for example, between prospective and retrospective studies or between
  RCTs and observational studies).
\item
  There is no information on effect-size measure (e.g.~SMD vs ratio).
\end{itemize}

\newpage

\subsection{SCORE (social science)}\label{score-social-science}

\textbf{Reference:} \textcite{tyner2026investigating}.

\textbf{Research question:} investigating the replicability of published
claims in the social and behavioural sciences.

\textbf{Data availability:} The SCORE replication package is available
at \url{https://osf.io/g5sny/}. The package includes replication data as
well as reviews of claims made in the included papers.

\textbf{Data description and source:} BEAR includes two SCORE outputs.
The source dataset of replications has 548 rows: 274 original claims and
274 matched replications. After BEAR processing and filtering, 267
matched replication rows are included. The source dataset for ``all
claims'' has 3,066 rows for claims from the set of replicated papers
(but no replications of these claims), where we select one statistic per
claim; after BEAR processing and filtering, 1,942 claim rows are
included.

\textbf{Data processing:} For the matched replication output, we used
the package's converted correlation scale where available for original
and replication statistics. For the set of all claims, effects are
heterogeneous across the source papers.

When several inputs are available for constructing a z-value in the set
of all claims, we prefer reported z, then reported t, then coefficient
divided by standard error, then signed square-root F for numerator df 1,
then a 95\% confidence interval with a point estimate, then a two-sided
p-value conversion. When the selected statistic only supplies a p-value
and no sign can be inferred, the z-value is unsigned. In the text for
all claims, it is typical to have multiple statistics backing up a
single claim, e.g.~``F(1,38) = 3.73, p = .033, partial eta-squared =
.16; F(1,39) = 7.28, p = .010, partial eta-squared = .16; F\textless1''.
We created a rule to pick one statistic per claim using agreement with
the SCORE classification as significant or non-significant, the
statistic's provenance, exactness, the available effect and standard
error, and text order.

\textbf{Study characteristics:}

\begin{itemize}
\item
  \textbf{Study ID:} we use the DOI/source-paper identifier in both
  outputs.
\item
  \textbf{Topic:} field of study (e.g.~``political science'') as
  classified by the authors.
\item
  The underlying study design is not systematically described.
\item
  Measures: SCORE ``claims'' dataset measures include regression
  coefficients, standardised mean differences, correlations, and
  variance-explained statistics; but the effect measure is unclassified
  for over 80\% of data. SCORE replications use correlations where
  available.
\end{itemize}

\newpage

\subsection{Lang 2025 (economics)}\label{lang-2025-economics}

\textbf{Reference:} \textcite{lang2025credible}.

\textbf{Research question:} How frequently are statistically significant
rejections of null hypotheses in recent empirical economics false
positives?

\textbf{Data collection:} Lang samples empirical articles from the top
five economics journals in 2021 and also articles previously collected
in Brodeur, Cook, and Heyes, who sampled articles published in 2015 and
2018 in 25 leading economics journals. (See also the ``Brodeur'' dataset
in BEAR for more details.) The sample consists mainly of natural,
laboratory, or field experiments, excluding articles focused on theory,
structural estimation, or econometric methods.

Lang re-extracted results from all articles rather than using the
hypothesis tests extracted by Brodeur, Cook, and Heyes. Research
assistants read the abstracts and results text and retained estimates
emphasised in the abstract or described by the authors as principal
results. When multiple specifications tested the same hypothesis, they
retained the authors' preferred specification or the specification whose
magnitude was used in subsequent analysis. Robustness checks and
heterogeneity analyses were excluded unless the heterogeneity analysis
was itself central to the article.

The source data contain 3,885 extracted hypothesis tests from 736
articles. Lang's main sample in the cited article restricts these to
rejected null hypotheses with \(1.96 \leq |t| < 10\), leaving 2,082
hypothesis tests from 663 articles. The dataset therefore does not
identify one focal hypothesis per article. An article can contribute
several hypothesis tests. However, only one preferred specification is
retained for each hypothesis. (Lang paper also reports analyses
restricted to one null hypothesis per article. The supplementary
material contains a separate vector with one t-statistic for each of the
736 articles, but this vector does not contain row or article
identifiers that allow it to be linked directly to the 3,885-test table.
We therefore do not use this in BEAR data.)

\textbf{Data processing:} We use the 3,885 results. We use Lang's
supplied, derounded t-statistic, derounded coefficient
(\texttt{coef\_dr} in the source data) and standard error
(\texttt{se\_dr}); the latter two are available in about 95\% of cases.
(Using either coefficients or t-statistics to calculate z-values leads
to practically identical results; we use the latter.) We calculate the
corresponding two-sided normal-approximation p-value. We retain some
additional columns in \texttt{Lang.rds}.

We also pull DOIs for source articles using bibliographic metadata
matching. These DOIs were obtained by us from Crossref, rather than
supplied by Lang; the saved assignments are preserved when reprocessing
the dataset. This is done automatically via Crossref. The matches are
very good in almost all cases. For several edge cases we used LLM review
to adjudicate.

\textbf{Study characteristics:}

\begin{itemize}
\item
  \textbf{Study ID:} \texttt{studyid} is \texttt{Lang\_paper\_} followed
  by Lang's \texttt{unique\_paperid}, but with a few minor corrections
  from us (see GitHub repository for details). Article DOIs are also
  available as an additional column in \texttt{Lang.rds}.
\item
  \textbf{Subset:} \texttt{subset} distinguishes observations drawn from
  the Brodeur--Cook--Heyes 2015/2018 source universe from Lang's 2021
  top-five additions.
\item
  \textbf{Method:} The \texttt{method} labels supplied with the Lang
  data mix study designs, identification strategies, and estimators.
  Large majority of estimates are one of ``RCT'', ``DID'', ``IV'',
  ``RD'', and ``OLS''. \texttt{Lang.rds} retains the original method
  labels. When constructing \texttt{BEAR.rds}, we collapse some label,
  e.g.~\texttt{RCT\ (DID)} and \texttt{RCT\ (IV)} become \texttt{RCT};
  \texttt{DID\ (matching)} becomes \texttt{DID}; and \texttt{RD\ (DID)}
  and \texttt{RD\ (IV)} become \texttt{RD}. For estimates labelled
  \texttt{OLS}, \texttt{matching}, \texttt{SYSTEM-GMM}, or
  \texttt{IV-DID\ (matching)}, we retain the estimates but set
  \texttt{method} to missing.
\item
  The dataset does not have data on type of measure to allow us to
  classify estimates as in some other datasets in BEAR (e.g.~risk risk
  vs standardised mean difference).
\end{itemize}

\newpage

\subsection{Szucs and Ioannidis 2017 (cognitive
neuroscience)}\label{szucs-and-ioannidis-2017-cognitive-neuroscience}

\textbf{Reference:} \textcite{szucs2017empirical}.

\textbf{Research question:} empirical assessment of published effect
sizes, statistical significance, and power in the literature on
cognitive neuroscience and psychology.

\textbf{Data availability:} The article and supporting data are
available from PLOS Biology at
\url{https://journals.plos.org/plosbiology/article?id=10.1371/journal.pbio.2000797}.

\textbf{Data description and source:} The source contains 26,841
text-mined t-test records from 3,801 papers in 18 journals (cognitive
neuroscience, psychology, and medical). The processed dataset preserves
all three fields; in the main BEAR dataset, we include only the
cognitive neuroscience subset (since the psychology and medical
literature are well covered by other datasets in BEAR).

\textbf{Notes:} The extraction targeted t-test records reported in
article texts and did not mine tables.

\textbf{Data processing:} We followed our default procedure for deriving
z-values from t-values.

\textbf{Study characteristics:}

\begin{itemize}
\item
  \textbf{Study ID:} to make look-ups of articles easier, we use
  \texttt{journal\_\textless{}journal\_code\textgreater{}\_article\_\textless{}article\_id\textgreater{}}.
\item
  \textbf{Topic:} we use the source \texttt{field} classification; note
  that \texttt{BEAR.rds} contains only cognitive neuroscience rows.
\item
  The source data does not classify study designs.
\item
  The original dataset labels all results as coming from a t-test, but
  there is no additional information on what was the underlying measure.
\end{itemize}

\newpage

\subsection{ClinicalTrials.gov
snapshot}\label{clinicaltrials.gov-snapshot}

Trials that reported results at ClinicalTrials.gov
\autocite{clinicaltrials2025}.

\textbf{Snapshot date:} 12 August 2025

\textbf{Data source:} This is a set of flat files of studies obtained
from the AACT downloads page
(\url{https://aact.ctti-clinicaltrials.org/downloads/}; over 2 GB of
data for relevant tables), with relevant columns extracted and joined by
us before BEAR processing.

\textbf{Data availability:} The terms and conditions for
ClinicalTrials.gov are at
\url{https://clinicaltrials.gov/about-site/terms-conditions}. We make
distributions of derived effect sizes available in accordance with the
terms and conditions. We only use publicly available data and do not
imply any endorsement of this work by the National Library of Medicine,
ClinicalTrials.gov or AACT.

\textbf{Data processing:} Processing was extensive. The full data table
(\texttt{BEAR\_data/clinicaltrialsgov.rds}) differs in some respects
from the main BEAR dataset. We briefly discuss these differences here.
Because of its rich structure, please refer to the data dictionary for
the full data table.

Briefly, BEAR targets primary outcome results for completed studies
using two data sources:

\begin{enumerate}
\def\labelenumi{(\arabic{enumi})}
\item
  Author-reported effect sizes
\item
  Effect sizes derived from raw outcome data (group means, counts),
  alongside source data
\end{enumerate}

Where both (1) and (2) are available, we prefer (1), but also make (2)
available in the full data table. For (2), we use standardised mean
differences for continuous outcomes and probit differences for count
outcomes.

When working with the full dataset, variable
\texttt{author\_raw\_overlap\_key} can be used to link the two different
calculations of the same effect; variable \texttt{include\_in\_bear}
identifies rows included after applying this rule preferring
author-reported estimates.

The main BEAR dataset keeps studies that (1) are randomised, (2) are
interventional, and (3) have fewer than 20 candidate effect rows (see
the description in this section).

The full ClinicalTrials.gov table uses only completed studies but does
not apply these three restrictions. It contains 151,313 rows from 27,446
studies; 46.3\% are author-reported and 53.7\% are derived from raw
data. The randomised interventional subset used in the main BEAR release
contains 60,470 rows from 23,060 studies.

Note that a single ClinicalTrials.gov study can contribute several
primary outcome IDs, and each primary outcome can contain many
measurements. The full data table records this as
\texttt{n\_effect\_rows\_per\_study} and
\texttt{n\_outcomes\_per\_study}. The median study contributes one
outcome ID and two rows to the full data table. Studies with 20 or more
candidate effect rows remain in the full data table but are excluded
from the main BEAR build.

In the full data table, some rows derived from raw data have
\texttt{direction\_unknown\ ==\ TRUE}, meaning that the
control/reference arm could not be reliably inferred, so the sign should
be treated as random. Users interested in controlled trials or signed
z-values should filter the results accordingly.

See \emph{Standard procedure for dealing with p-values and confidence
intervals} for how we created and chose the z-values based on p-values
and confidence intervals. Briefly, we calculated z-values using two
separate procedures (based on p-values and confidence intervals) and, if
both were available, we made a judgement on which one was more
appropriate using a predefined rule. ClinicalTrials.gov reports the
number of sides for confidence intervals but not for p-values, so we
treated p-values as two-sided. We dropped rows only when \texttt{z} was
missing. Notably, infinite z-values can arise when \texttt{p\_value} is
recorded as exactly 0, and these are included because they are not
missing.

\textbf{Study characteristics:}

\begin{itemize}
\item
  \textbf{Study ID:} we use the ClinicalTrials.gov trial identifier
  (\texttt{nct\_id}).
\item
  \texttt{subset} records trial phase.
\item
  We include only randomised interventional studies, including parallel,
  crossover and factorial designs.
\item
  About half of all effects (by which we mean both reported and our
  derived estimates) are standardised mean differences, mean
  differences, or probit differences. Another 14\% are hazard, odds,
  risk, geometric, or rate ratios; 5\% have the other category (in
  \texttt{BEAR.rds} this is coded as ``other''), and 20\% have no
  measure.
\end{itemize}

\newpage

\subsection{EU Clinical Trials Register (EU
CTR)}\label{eu-clinical-trials-register-eu-ctr}

Trials that reported results in the EU Clinical Trials Register,
\url{https://www.clinicaltrialsregister.eu/} \autocite{euctr2025}.

\textbf{Snapshot date:} 18 January 2026

\textbf{Data collection:} We used the R package \texttt{ctrdata}
\autocite{herold2025ctrdata} to create a snapshot of EU CTR
(\texttt{collection\ =\ "euctr"}); we downloaded only trial records with
results. We stored the downloaded records in a local SQLite database and
extracted a fixed set of protocol and results fields plus several fields
computed by \texttt{ctrdata} (trial phase, sample size, and a ``first
primary endpoint'' p-value/size derived by the package). We use only the
EU CTR portion of this extraction in BEAR.

\textbf{Data availability:} EUCTR is publicly accessible. The EUCTR
legal notice notes that publication on EUCTR does not constitute an
endorsement of the information reported.

\textbf{Data processing:} Processing was extensive. We use only primary
endpoints. We default missing CI levels to 95\% and where truncation of
p-values was not stated we assume equality (the latter is the case for
about 15\% of data). Where the data do not report whether a test is one-
or two-sided, we assume two-sided. The rest of the derivation of
z-values follows the \emph{Standard procedure for dealing with p-values
and confidence intervals}. We set \texttt{year} to the year of entry
since completion dates were unavailable. We dropped rows only when
\texttt{z} was missing; infinite z-values are included. We harmonise
phase labels to follow the same naming convention as ClinicalTrials.gov.

\textbf{Study characteristics:}

\begin{itemize}
\item
  \textbf{Sample size:} \texttt{ss} is the extracted endpoint sample
  size (\texttt{n}).
\item
  \textbf{Study ID:} we use the EudraCT trial identifier (\texttt{id}).
\item
  \texttt{subset} records trial phase.
\item
  The registry extract does not provide a row-level randomisation
  indicator, so BEAR does not classify the trial design.
\item
  Effect measures are derived from reported estimand labels; about half
  of these are mean differences, hazard ratios, odds ratios, risk
  differences, or risk ratios. Another one third have other category (in
  \texttt{BEAR.rds} this is coded as ``other''), and 23\% have no
  measure.
\item
  \texttt{subset} records registry phase.
\end{itemize}

\newpage

\subsection{Barnett and Wren 2019
(PubMed/Medline)}\label{barnett-and-wren-2019-pubmedmedline}

\textbf{Reference:} \textcite{barnett2019examination}.

\textbf{Research question:} Bias for statistical significance in health
and medical journals

\textbf{Data availability:} The dataset \texttt{Georgescu.Wren.RData}
may be downloaded from
\url{https://github.com/agbarnett/intervals/tree/master/data} or
\url{https://github.com/jdwren/ASEC}. The repositories give no clear
licence, but the BMJ Open article is CC BY 4.0.

\textbf{Data description and source:} The dataset is a collection of
confidence intervals for ratio estimates from Medline papers published
between 1976 and 2019. Authors scraped (via regular expressions,
followed by an independent check using another data mining algorithm,
manual checks for 10,000 intervals) 968,000 CIs from abstracts and
350,000 from full texts.

\textbf{Notes:} These are ``ratio estimates'' such as odds ratios,
hazard ratios and risk ratios. Note that binary outcomes tend to contain
(much) less information than continuous outcomes.

\textbf{Data processing:} We exclude a small fraction of data. First, we
restrict the data to 95\% confidence intervals: where \texttt{ci.level}
was missing we assume it was 0.95; among intervals with a known
\texttt{ci.level}, about 0.3\% were not 95\% and are dropped. We drop
intervals with non-positive width
(i.e.~\texttt{lower\ \textgreater{}=\ upper}). We use the log scale (we
replace zero or negative lower bounds with a small positive constant;
this affects about 0.1\% of the sample) and back out an approximate
standard error and point estimate under the usual normal approximation
for a 95\% CI, setting the point estimate to the midpoint of the
interval.

For the main version of BEAR dataset, \texttt{BEAR.rds}, 50,000 studies
are selected at random, with one row per study, in order to keep the
file size manageable.

\textbf{Study characteristics:}

\begin{itemize}
\item
  \textbf{Study ID:} we use the PubMed identifier (\texttt{pubmed}).
\item
  There is no information on study design (i.e.~we cannot distinguish,
  for example, between prospective and retrospective studies or between
  RCTs and observational studies).
\item
  Effects are ratios on a logarithmic scale, without distinguishing
  odds, risks, hazards etc.
\end{itemize}

\newpage

\subsection{Open Science Collaboration 2015
(psychology)}\label{open-science-collaboration-2015-psychology}

\textbf{Reference:}
\textcite{opensciencecollaborationEstimatingReproducibilityPsychological2015a}.

\textbf{Research question:} replications of a quasi-random sample of 100
experiments in psychology.

\textbf{Data collection:} The authors make a comprehensive dataset of
results available with their paper. BEAR includes only the replications
because all original studies in that set have \(p < 0.05\).

\textbf{Data availability:} The data file \texttt{rpp\_data.csv} is
available in the paper's repository,
\url{https://osf.io/ytpuq/overview}, under a CC0 1.0 Universal licence.

\textbf{Data processing:} We derived z-values from the replication
p-values assuming two-sided p-values (we use the p-value column for the
replication, not the p-value for the original study). We applied the
sign of the replication correlation, so z-values are signed where that
correlation is available. We treated p-values recorded as exactly 0 as
truncated and recorded this as a lower bound on z (z-operator
\texttt{\textgreater{}}); otherwise we treated p-values as exact
(z-operator \texttt{=}). We retained the replication sample size as
\texttt{ss}.

In \texttt{BEAR.rds}, the measure is recorded as ``correlation'', based
on the fact that is the column from original dataset that we use.

\textbf{Study characteristics:}

\begin{itemize}
\item
  BEAR treats the replication studies as RCTs.
\item
  The retained effect measure is the replication correlation.
\end{itemize}

\newpage

\section{Standard procedure for dealing with p-values and confidence
intervals}\label{standard-procedure-for-dealing-with-p-values-and-confidence-intervals}

Most datasets in BEAR come with pre-computed effect sizes and standard
errors. Some, especially databases of clinical trials and large sets of
scraped studies, report p-values. Moreover, reporting in clinical trial
databases is done by investigators, and the data-entry fields are not
standardised, requiring more extensive clean-up and a procedure for
deriving z-values. Lastly, in some cases datasets contain confidence
intervals and effect sizes, but z- and p-values may be missing.

\subsection{Data cleaning and default
values}\label{data-cleaning-and-default-values}

For confidence interval levels, we set implausible values to missing
(e.g.~values above 100 or below 0) and standardise proportions and
percentages: investigators sometimes report 90 and sometimes 0.9 for a
90\% confidence interval. In EU CTR, confidence levels can be missing,
and we assume 95\% by default. Some data state whether CIs or p-values
are one- or two-sided.

Similarly, when working with p-values we assume a two-sided test when no
additional information is stated.

\subsection{Confidence intervals}\label{confidence-intervals}

When working with CIs, we analyse ratio measures on the log scale when
the lower bound, upper bound, and point estimate are all strictly
positive; we keep non-ratio measures on the raw scale.

To derive SEs from confidence intervals, we first compute the critical
value. To avoid undefined critical values (e.g.~at 100\% or extremely
close to it), we bound the implied \(\alpha\) at a small value. We then
compute two one-sided standard error estimates from the upper and lower
bounds and use their average when both are finite.

We compute a simple symmetry diagnostic for the confidence interval on
the chosen scale, defined as the smaller of the two half-widths divided
by the larger. We treat intervals as ``symmetric enough'' only if this
ratio exceeds 0.8. This arbitrary threshold affects when we use the
CI-derived z-statistic; otherwise, we use a p-derived z-statistic where
one is available.

Where additional information is available, we look for potentially
non-Wald intervals by matching text in the parameter-type column. We
search for keywords such as median, Hodges, posterior/Bayes, exact,
Fieller, bootstrap, and permutation. For flagged estimates, we prefer to
derive the z-value from the p-value.

\subsection{p-values}\label{p-values}

When a source reports t-statistics and degrees of freedom, we first
derive the two-sided p-value from the t distribution and then derive the
z-value (retaining the sign of the original t-statistic, if available).

For clinical trials databases (ClinicalTrials.gov and EU CTR), we
sometimes find ambiguity in how p-values are recorded. Some significant
results report p-values close to 1, suggesting that authors report the
cumulative distribution, \(\Phi(z)\). Therefore, for each row we also
calculate \(2\cdot\min\{p,1-p\}\). When a CI-derived z-value is
available, we choose the p-value that gives the closest z-value.

To avoid dropping rows, we retain unsigned z-statistics when the sign
cannot be determined. If the sign of the effect is known, we apply it to
the p-derived z; otherwise, we use its absolute value.

\subsection{Choice of a single z-value when multiple candidates are
available}\label{choice-of-a-single-z-value-when-multiple-candidates-are-available}

In situations where \(z\) can be derived from both CIs and p-values, we
choose a single ``best'' z-statistic per row as follows. We prefer the
CI-derived \(z\) when it looks like a plausible Wald interval on the
chosen scale, meaning: (1) finite \(z\) and standard error, (2) positive
standard error, (3) no non-Wald flag from a keyword search of the
measure label, and (4) either missing symmetry information or symmetry
above 0.8. Otherwise, we use the p-derived z.

\subsection{Truncation and infinite
z-values}\label{truncation-and-infinite-z-values}

For reported two-sided p-value bounds, \(p<c\) gives a lower bound on
\(|z|\). BEAR records this inequality using \texttt{z\_operator}. The
operator applies to the magnitude, including when the stored z-value is
negative. For example, a stored value of \(-1.96\) with operator
\texttt{\textgreater{}} means a negative direction and magnitude greater
than 1.96, not \(z>-1.96\).

Exact zero p-values can produce infinite z-values. We retain infinities
in BEAR source data, although signal-to-noise models truncate these
infinite values.
